% AI-produced by GPT-6 Astra (OpenAI, ChatGPT Pro mode).
% Reconstructed on 2026-11-20 from the complete accepted in-session argument,
% including the final collar and relative-normal-bundle refinements.
% The historical TeX bytes were lost; this is a documented new reconstruction.
\Needspace{10\baselineskip}
\subsection{A.146. Vanishing of the component obstruction for the three-torus}

\begin{saresult}{A.146}{Proved}{The formal ordinary-concordance obstruction vanishes on the three-torus}
For the standard real PL torus $T^3$ the homomorphism of A.143 is zero:
\[
             c_{T^3}:\pi_0P_{\mathrm{TOP}}(T^3)
                         \longrightarrow H^2(T^3;\mathbb Z/2),
             \qquad c_{T^3}=0.
 \tag{146.1}
\]
This is a computation of the obstruction homomorphism. It does not
assert that every TOP concordance is PL-isotopic, or TOP-isotopic,
to a PL concordance.
\end{saresult}

\begin{proof}
Put $X=T^3\times I$. Start with any bottom-fixed TOP concordance.
Its top restriction is homotopic to the identity by projection of
the concordance to $T^3$. The endpoint normalization in A.143 puts
it in the same free-top component as a map with PL top restriction
$a$, still homotopic to the identity. Waldhausen's theorem for the
sufficiently large irreducible three-manifold $T^3$ makes $a$
TOP-isotopic to the identity; the exact original locator is
\cite[Theorem 7.1, printed pp.~80--81]{Waldhausen1968}.
The injectivity on components of the actual joint-family
PL-to-TOP automorphism comparison in A.134 then supplies a jointly
PL path from $a$ to the identity. This last step avoids extracting
an unproved joint-PL convention from Waldhausen's isotopy wording.
Its PL trace corrects the top. We have therefore reduced the value
of $c_{T^3}$ to that of a homeomorphism
\[
                  F:X\longrightarrow X,\qquad
                         F|_{\partial X}=1.
 \tag{146.2}
\]
All endpoint corrections are the trace corrections already justified
in A.143; PL traces have zero formal obstruction.

First make $F$ the identity on collars of both time ends, by the
boundary compression isotopy (145.5). Lift its first coordinate to
$\mathbb R^3\times I$, choosing the lift equal to $x$ at the bottom.
The lift is equivariant for $\mathbb Z^3$ because the bottom
inclusion induces an isomorphism on fundamental groups and $F$
fixes that bottom. At the top its first coordinate is $x+z$ for
one $z\in\mathbb Z^3$. Choose a PL function
$\rho:I\to I$ equal to zero near $0$ and one near $1$, and compose
with the PL homeomorphism
\[
        L_z([x],s)=([x-\rho(s)z],s).
 \tag{146.3}
\]
It is the identity near both ends: at the top the translation by
the integral vector $z$ is the identity of the torus. It is PL
because it is induced by a finite piecewise affine map in the
standard torus product coordinates. Its component has zero value
under $c_{T^3}$. Thus replacing $F$ by $L_zF$ leaves the obstruction
unchanged and makes the chosen top lift equal to $x$.

Write the resulting map, still denoted $F$, as
\[
                 F([x],s)=([x+w(x,s)],v(x,s)),
 \tag{146.4}
\]
where $w$ is $\mathbb Z^3$-periodic and continuous. On collars of
$s=0,1$ we have $w=0$ and $v=s$. Hence
\[
 k_t([x],s)=
   \bigl([x+t\,w(x,s)],\, (1-t)s+t\,v(x,s)\bigr),
          \qquad 0\le t\le1,
 \tag{146.5}
\]
is a homotopy of maps from the identity to $F$, fixed on those
collars. The maps $k_t$ are not asserted to be homeomorphisms.
Reparameterize $t$ by a continuous PL function constant near its
two endpoints. The homotopy then equals the identity for $t$
near zero and $F$ for $t$ near one, while it remains fixed near
the full spatial boundary.

Form the labeled mapping cylinder
\[
 M_F=(X\times I)\cup_F X.
\]
Its top copy of $X$ is labeled by $y=F(x)$ for the cylinder point
$(x,1)$. Define
\[
 U:M_F\longrightarrow Z=X\times I,\qquad
             U([x,t])=(k_t(x),t).
 \tag{146.6}
\]
The formula is well-defined at the top precisely because $k_1=F$;
there it is the identity in the top label $y$. It is also the
identity in the bottom and side labels. The preceding collar
choices make it literally product in the lower, upper, and side
boundary coordinates, including their corner intersections.
Both bottom inclusions are homotopy equivalences, and $U$
restricts to their common identity. Consequently $U$ is a
homotopy equivalence. This constructs an actual relative
structure, rather than postulating a normal invariant for $F$.

The target $Z=T^3\times I^2$ is a compact five-manifold. A homotopy
equivalence of this full boundary pair determines its ordinary
relative TOP normal invariant
\[
               \nu(U)\in[Z,\partial Z;G/\mathrm{TOP},*].
 \tag{146.7}
\]
Use the source stable TOP normal bundle and the induced
fiber-homotopy equivalence of Spivak normal fibrations to define
the normal data. The product boundary labels specify the relative
identification there. Since $U$ is a homotopy equivalence, its
surgery obstruction is zero. Also
$\operatorname{Wh}(\mathbb Z^3)=0$, so there is no ambiguity between
simple and homotopy decorations in this application.

We now use the assembly calculation in its correct dimension.
Definition 18.2, printed pp.~197--198, of
\cite{Ranicki1992TopologicalManifolds} states the surgery sequence
for compact manifolds with boundary of dimension at least five.
The relative-version paragraph on printed p.~198, Proposition 18.3,
and Theorem 18.5, printed pp.~198--201, identify its normal-invariant
map with assembly. These inputs apply to the five-dimensional
$Z$, not to the three-dimensional base. Poincare--Lefschetz
duality identifies the normal-invariant group with
\[
 [Z,\partial Z;G/\mathrm{TOP},*]
           \cong H_5(Z;\mathbb L\langle1\rangle)
           \cong H_5(T^3;\mathbb L\langle1\rangle).
 \tag{146.8}
\]
Here $\mathbb L\langle1\rangle$ denotes the one-connective quadratic
surgery spectrum.

The same original, Example 24.16, printed p.~278, gives the
assembly isomorphism for free abelian groups in degrees at least
their rank, and the vanishing of their Whitehead groups. In
particular,
\[
 H_5(T^3;\mathbb L\langle0\rangle)
       \xrightarrow{\ \cong\ }L_5(\mathbb Z[\mathbb Z^3]).
 \tag{146.9}
\]
Passing from the one-connective to the zero-connective spectrum
does not alter the degree-five group here: the intervening
ordinary-homology groups are $H_5(T^3;\mathbb Z)$ and
$H_4(T^3;\mathbb Z)$, both zero. Thus (146.8)--(146.9) make the
relative surgery obstruction homomorphism injective. Since
$\theta(\nu(U))=0$, we obtain
\[
                             \nu(U)=0.
 \tag{146.10}
\]

We spell out why (146.10) kills the \emph{specified relative}
PL obstruction. Forget the spherical trivialization in the
relative normal invariant and retain the stable TOP normal
reduction, with its fixed boundary comparison. The normal
invariant compares the source normal bundle, transported across
$U$, with the standard PL normal reduction on $Z$. Its zero class
gives the corresponding relative homotopy of TOP reductions to
that standard reduction. Pulling back the PL lift along this
homotopy gives a stable PL normal reduction on $M_F$ agreeing
with its prescribed boundary reductions. This is also seen by
applying the natural universal PL-obstruction map to the
relative normal invariant: the resulting class is exactly the
relative Kirby--Siebenmann class of the labeled source, since
the target's standard PL reduction has zero class. The relative
collar identifications in (146.6) ensure that no boundary
difference is discarded.

The stable tangent and normal formulations give the same
mod-two obstruction, by the Whitney-sum formula and the stable
trivial tangent-plus-normal bundle. Therefore
\[
                  k(M_F,\partial M_F)=0.
 \tag{146.11}
\]
By the mapping-cylinder definition (G.5) and the interval
suspension isomorphism this says
$c_{\mathrm{PL}}(F)=0$. The preliminary trace and PL corrections
left $c_{T^3}$ unchanged, so (146.1) follows for the original
ordinary concordance.

Only relative normal invariants in dimension five and the
explicit torus assembly isomorphism were used. A stable
normal-bundle reduction was not converted into an isotopy of
four-dimensional homeomorphisms. Such a conversion would be an
additional assertion and is not part of this theorem.
\end{proof}

